% ============================================================================
%  SECTION 1 — Biological motivation & fundamental observations
%  Order: motivation plots first; constituents introduced only at the end.
% ============================================================================
\begin{frame}[t]{}
  \vfill\begin{center}{\Huge\color{YaleBlue} The biology}
  \end{center}\vfill
\end{frame}

\begin{frame}[t]{Living Materials Are Never ``Finished''}
  A steel beam is made once and wears out. A living artery is
  \textbf{continuously torn down and rebuilt}: cells deposit and degrade
  structural proteins while sensing the loads they carry.
  \begin{itemize}
    \item \textbf{Growth} --- change in \emph{mass}.
    \item \textbf{Remodeling} --- change in \emph{microstructure} (fiber
    orientation, natural length, composition).
  \end{itemize}
  \vfill
  \begin{center}\color{AccentRed}
  Will a hypertensive artery thicken and stabilise, or an aneurysm dilate to
  rupture?
  \end{center}
\end{frame}

\begin{frame}[t]{Cells Actively Hold a Mechanical Set-Point}
  Cultured cells generate a steady \textbf{tension} and \emph{restore} it after
  each perturbation --- direct evidence that living tissue targets a mechanical
  state rather than a fixed shape.
  \bigfigw{biol_tensional_homeostasis.png}{0.72\linewidth}{Force generated by a
  contractile cell line vs.\ control, recovering after repeated unloading.}
  \srccite{Tensional homeostasis; Brown et al., J Cell Physiol (1998)}
\end{frame}

\begin{frame}[t]{Fundamental Observation: A Stress Set-Point}
  Wolinsky \& Glagov (1967): Cauchy stress in the aortic media is nearly
  \textbf{constant across mammals} spanning orders of magnitude in body mass ---
  tissue-level evidence of a mechanical set-point.
  \bigfigw{motiv_wolinsky_glagov.png}{0.6\linewidth}{Wall tension vs.\ body mass
  across species.}
  \srccite{Wolinsky \& Glagov, Circ Res (1967)}
\end{frame}

\begin{frame}[t]{Cells Chase a Mechanical Set-Point}
  Cells regulate a \textbf{target (homeostatic) stress}: above it they build more,
  below it removal wins --- a negative feedback that thickens a hypertensive wall
  until stress returns to normal.
  \vfill
  \begin{center}\gnrloop[0.86]\end{center}
  \vfill
  \srccite{Humphrey, J Elasticity (2021)}
\end{frame}

\begin{frame}[t]{Two Canonical Scenarios We Will Simulate}
  \begin{columns}[T]
    \begin{column}{0.5\textwidth}
      \textbf{\color{YaleBlue}Hypertension} --- sustained pressure step. Expect
      \emph{adaptation}: the wall thickens, stress returns to set-point.
    \end{column}
    \begin{column}{0.5\textwidth}
      \textbf{\color{YaleBlue}Aneurysm} --- irreversible elastin loss. Geometry can
      feed \emph{positively} back on stress: stabilise larger, \emph{or} dilate
      without bound.
    \end{column}
  \end{columns}
  \vfill
  \begin{center}\large\color{AccentRed}
  \emph{When does a tissue adapt, and when does it run away?}
  \end{center}
\end{frame}

% ---- constituents introduced only at the END of the biology section --------
\begin{frame}[t]{The Cast of Characters in an Arterial Wall}
  \footnotesize
  \begin{center}
  \renewcommand{\arraystretch}{1.35}
  \begin{tabular}{@{}p{0.135\textwidth}p{0.235\textwidth}p{0.115\textwidth}p{0.115\textwidth}p{0.30\textwidth}@{}}
    \toprule
    \textbf{Constituent} & \textbf{Mechanical role} & \textbf{Mass frac.\ $\phi$} &
      \textbf{Dep.\ stretch $G$} & \textbf{Material model \& turnover $1/k_d$}\\
    \midrule
    \textbf{\color{YaleBlue}Elastin} & soft; bears load at low stretch & $0.30$ & $1.40$ &
      neo-Hookean ($c\approx90$\,kPa); \textbf{none} (decades) \\
    \textbf{\color{YaleBlue}Collagen} & stiff, crimped; guards against rupture & $0.35$ & $1.08$ &
      Fung ($c_1\approx250$\,kPa, $c_2\approx8$); \textbf{fast} ($\sim$20\,d) \\
    \textbf{\color{YaleBlue}Smooth muscle} & passive load + active contraction; the sensor cell &
      $0.35$ & $1.20$ & Fung ($c_1\approx120$\,kPa, $c_2\approx3$); fast ($\sim$20\,d) \\
    \bottomrule
  \end{tabular}
  \end{center}
  \vfill
  \begin{itemize}\footnotesize
    \item \textbf{Elastin cannot be renewed} --- its loss triggers \emph{aneurysm}.
    \item \textbf{Collagen \& SMC turn over constantly} --- so the wall can
    \emph{adapt}, or \emph{run away}.
  \end{itemize}
  \begin{center}\scriptsize\color{BoxGray}
  Teaching-grade values for a large elastic artery, in the range used by Humphrey,
  Cyron, Latorre and co-workers (order-of-magnitude, not a fit to one animal).
  \end{center}
  \srccite{Humphrey, Eberth, Dye, Gleason, J Biomech (2009)}
\end{frame}
